% ===========================================================
% Diseñados para amar — Preámbulo
% Formato: A5, dos caras, libro de bolsillo
% Compilar con: lualatex o xelatex
% ===========================================================

% --- Geometría de página ---
\usepackage[
  a5paper,
  inner=20mm,
  outer=15mm,
  top=18mm,
  bottom=20mm,
  bindingoffset=3mm
]{geometry}

% --- Idioma ---
\usepackage[spanish]{babel}

% --- Tipografía (requiere LuaLaTeX o XeLaTeX) ---
\usepackage{fontspec}
\setmainfont{EB Garamond}[
  Numbers=OldStyle,
  Ligatures=TeX
]
% Fallback si no está EB Garamond:
% \setmainfont{Libertinus Serif}

% --- Interlineado ---
\usepackage{setspace}
\setstretch{1.3}

% --- Mejoras tipográficas ---
\usepackage{microtype}

% --- Imágenes ---
\usepackage{graphicx}
\graphicspath{{../img/}}

% --- TikZ para elementos decorativos ---
\usepackage{tikz}
\usetikzlibrary{calc, decorations, patterns, fadings}

% Colores del libro — paleta azul noche + oro cálido
\definecolor{noche}{HTML}{1B2838}       % Deep navy / midnight
\definecolor{oro}{HTML}{D4A843}          % Warm gold
\definecolor{negro}{HTML}{1A1A1A}       % Near-black for text
\definecolor{gris}{HTML}{6B7B8D}        % Slate grey for accents
\definecolor{portada}{HTML}{0C2233}     % Deep petrol blue for cover

% --- Tipografía sans-serif para portada ---
\setsansfont{Calibri}[
  Ligatures=TeX
]

% --- Comando: vesica piscis (geometría sagrada) ---
% Dos círculos solapados que crean una mandorla
% Uso: \vesicapiscis[opacidad]{x}{y}{radio}
\newcommand{\vesicapiscis}[4][0.15]{%
  \begin{scope}[shift={(#2,#3)}, opacity=#1]
    \draw[noche, line width=0.5pt] (-0.5*#4,0) circle (#4);
    \draw[oro, line width=0.5pt]   (0.5*#4,0) circle (#4);
  \end{scope}
}

% --- Comando: retícula de vesica piscis ---
\newcommand{\fondovesica}[1][0.06]{%
  \foreach \row in {0,1,...,8} {
    \foreach \col in {0,1,...,5} {
      \pgfmathsetmacro{\xpos}{-1.5 + \col * 3.0}
      \pgfmathsetmacro{\ypos}{-14 + \row * 3.0}
      \vesicapiscis[#1]{\xpos}{\ypos}{1.1}
    }
  }
}

% --- Comando: línea decorativa dorada ---
\newcommand{\lineadecorativa}[1][3cm]{%
  \begin{tikzpicture}
    \draw[oro, line width=0.6pt] (0,0) -- (#1,0);
    \fill[oro] (-0.08,0) circle (1.2pt);
    \fill[oro] (#1+0.08,0) circle (1.2pt);
  \end{tikzpicture}%
}

% --- Ajuste de headheight ---
\setlength{\headheight}{12.5pt}
\addtolength{\topmargin}{-0.5pt}

% --- Encabezados y pies ---
\usepackage{fancyhdr}
\pagestyle{fancy}
\fancyhf{}
\fancyhead[LE]{\small\itshape Diseñados para amar}
\fancyhead[RO]{\small\itshape\leftmark}
\fancyfoot[C]{\thepage}
\renewcommand{\headrulewidth}{0.4pt}
\renewcommand{\footrulewidth}{0pt}

% Estilo para páginas de inicio de capítulo
\fancypagestyle{plain}{
  \fancyhf{}
  \fancyfoot[C]{\thepage}
  \renewcommand{\headrulewidth}{0pt}
}

% --- Títulos de capítulo ---
\usepackage{titlesec}

% --- Formato de \part con página decorada ---
% Fondo azul-petróleo semitransparente (eco de la portada)
\newcommand{\partpage}[3]{%
  % #1 = número de parte (I, II…)
  % #2 = título de la parte
  % #3 = subtítulo de la parte
  \cleardoublepage
  \stepcounter{part}%
  \thispagestyle{empty}
  \begin{tikzpicture}[remember picture, overlay]
    % --- Fondo azul-petróleo con transparencia ---
    \fill[portada, opacity=0.70]
      (current page.north west) rectangle (current page.south east);
    % --- Vesica piscis sutil ---
    \begin{scope}[shift={(current page.center)}, yshift=0.5cm]
      \draw[white, line width=2.5pt, opacity=0.15] (-1.8,0) circle (4cm);
      \draw[oro,   line width=2.5pt, opacity=0.15] (1.8,0)  circle (4cm);
    \end{scope}
    % --- Marco fino dorado ---
    \draw[oro, opacity=0.25, line width=0.4pt]
      ([xshift=10mm, yshift=-10mm]current page.north west) rectangle
      ([xshift=-8mm, yshift=10mm]current page.south east);
  \end{tikzpicture}
  \vspace*{0.25\textheight}
  \begin{center}
    {\color{oro}\lineadecorativa[2cm]}\\[1.2em]
    {\sffamily\fontsize{14}{16}\selectfont\color{oro!80!white}\scshape Parte #1}\\[1.5em]
    {\sffamily\fontsize{24}{28}\selectfont\bfseries\color{white}#2}\\[1em]
    {\color{oro}\lineadecorativa[2cm]}\\[1.2em]
    {\sffamily\fontsize{12}{15}\selectfont\itshape\color{oro!80!white}#3}
  \end{center}
  \vfill
  \clearpage
}

\titleformat{\chapter}[display]
  {\normalfont\Large\bfseries}
  {\chaptertitlename\ \thechapter}
  {10pt}
  {\LARGE}

\titleformat{\section}
  {\normalfont\large\bfseries}
  {\thesection}
  {1em}
  {}

\titleformat{\subsection}
  {\normalfont\normalsize\bfseries\itshape}
  {\thesubsection}
  {1em}
  {}

% --- Notas al pie ---
\usepackage[bottom, hang]{footmisc}
\setlength{\footnotemargin}{1.2em}
\renewcommand{\footnoterule}{%
  \kern-3pt
  \hrule width 0.33\textwidth height 0.4pt
  \kern 2.6pt
}
% Numeración por capítulo
\usepackage{chngcntr}
\counterwithin*{footnote}{chapter}

% --- Epígrafes ---
\usepackage{epigraph}
\setlength{\epigraphwidth}{0.75\textwidth}
\renewcommand{\epigraphflush}{flushright}
\renewcommand{\sourceflush}{flushright}

% --- Listas personalizadas ---
\usepackage{enumitem}

% --- Citas en bloque ---
\usepackage{csquotes}
\SetBlockEnvironment{quotation}

% --- Bibliografía ---
\usepackage[
  backend=biber,
  style=verbose-note,
  language=spanish
]{biblatex}
% \addbibresource{bibliografia.bib}

% --- Enlaces (último paquete) ---
\usepackage[
  hidelinks,
  pdfauthor={Jesús Torrecilla Pinero},
  pdftitle={Diseñados para amar},
  pdfsubject={Antropología del matrimonio y la familia},
  pdfkeywords={matrimonio, familia, amor, antropología}
]{hyperref}

% --- Comando para sección "Para caminar juntos" ---
\newcommand{\paracaminar}{%
  \vspace{1.5em}
  \noindent\rule{0.3\textwidth}{0.4pt}
  \vspace{0.5em}
  \subsection*{Para caminar juntos}
}

% --- Viñeta de fin de capítulo ---
% Uso: \vinetacapitulo{ch01} (sin extensión)
\newcommand{\vinetacapitulo}[1]{%
  \vspace{1.5em}
  \begin{center}
    \includegraphics[width=0.5\textwidth]{#1}%
  \end{center}
}
